\documentclass[a4paper,10pt]{article}
\usepackage[empty]{fullpage}
\usepackage{titlesec}
\usepackage[usenames,dvipsnames]{color}
\usepackage{enumitem}
\usepackage[hidelinks]{hyperref}
\usepackage{fancyhdr}

\pagestyle{fancy}
\fancyhf{}
\renewcommand{\headrulewidth}{0pt}
\renewcommand{\footrulewidth}{0pt}

\addtolength{\oddsidemargin}{-0.375in}
\addtolength{\evensidemargin}{-0.375in}
\addtolength{\textwidth}{0.75in}
\addtolength{\topmargin}{-0.5in}
\addtolength{\textheight}{1.25in}

\titleformat{\section}{\vspace{-4pt}\large\bfseries\scshape\raggedright}{}{0em}{}[\color{black}\titlerule\vspace{-5pt}]

\newcommand{\resumeSubheading}[4]{
  \item
    \begin{tabular*}{0.97\textwidth}{l@{\extracolsep{\fill}}r}
      \textbf{#1} & \textit{#2} \\
      \textit{\small#3} & \textit{\small#4} \\
    \end{tabular*}
}

\newcommand{\resumeItem}[2]{\item \small{\textbf{#1}{: #2 \vspace{-2pt}}}}
\newcommand{\resumeSubHeadingListStart}{\begin{itemize}[leftmargin=*]}
\newcommand{\resumeSubHeadingListEnd}{\end{itemize}}
\newcommand{\resumeItemListStart}{\begin{itemize}}
\newcommand{\resumeItemListEnd}{\end{itemize}}

\begin{document}

% Heading
\begin{center}
    \textbf{\Huge Praveen Rangavajhula} \\[2pt]
    \href{mailto:praveen.rangavajhula@uga.edu}{praveen.rangavajhula@uga.edu} $|$
    \href{https://www.linkedin.com/in/praveen-rangavajhula/}{LinkedIn} $|$
    \href{https://github.com/Praveen-Rangavajhula}{GitHub} $|$
    Mobile: +1-706-765-9109
\end{center}

% Profile summary
\section{Profile Summary}
First-year PhD student in Computer Science at the University of Georgia with research interests in machine learning, predictive video modeling, longitudinal machine learning, optimization, learning dynamics, and interpretability. Experience includes V-JEPA-based latent future prediction, cognitive outcome modeling across longitudinal cohorts, automatic differentiation systems, privacy-preserving optimization, and federated learning.

% Education
\section{Education}
\resumeSubHeadingListStart
    \resumeSubheading
        {University of Georgia (UGA)}{Athens, GA}
        {Ph.D. in Computer Science}{Jan 2026 - Present}
    \resumeSubheading
        {University of Georgia (UGA)}{Athens, GA}
        {M.S. in Computer Science (GPA: 3.91)}{Aug 2023 - Dec 2025}
    \resumeSubheading
        {BITS Pilani}{Pilani, India}
        {B.E. in Mechanical Engineering}{Aug 2019 - July 2023}
\resumeSubHeadingListEnd

% Research experience
\section{Research Experience}
\resumeSubHeadingListStart

    \resumeSubheading
        {Graduate Research Assistant}{University of Georgia}
        {Longitudinal Machine Learning for Cognitive Development using FFCWS and MIDUS}{Feb 2026 - Present}
        \resumeItemListStart
            \resumeItem{End-to-End ML Pipeline}
                {Owned the complete machine learning pipeline for predicting cognitive outcomes from longitudinal Future of Families and Child Wellbeing Study (FFCWS) data, including data ingestion, cross-wave alignment, preprocessing, train-only imputation, model implementation, evaluation, analysis, and visualization for 1,480 children across ages 3, 5, and 9.}
            \resumeItem{Model Development and Benchmarking}
                {Implemented and benchmarked nine model families: standard and feature-engineered linear regression, support vector regression, Random Forest, XGBoost, histogram gradient boosting, an attention network, TabICLv2, and a custom Neural CDE. Evaluated the models using five repeated child-level 80/20 train-test splits and 344 to 796 time-valid predictors.}
            \resumeItem{Empirical Findings}
                {Found that expanding from theory-driven features to hundreds of predictors increased TabICLv2's average held-out $R^2$ from 0.307 to 0.324, indicating limited incremental predictive value from the expanded feature space.}
            \resumeItem{Cross-Dataset Extension}
                {Extending the end-to-end pipeline to the Midlife in the United States (MIDUS) longitudinal study to examine whether the findings generalize to an independent cohort.}
        \resumeItemListEnd

    \resumeSubheading
        {Master's Researcher}{University of Georgia}
        {Advised by Prof. John A. Miller}{Aug 2024 - Dec 2025}
        \resumeItemListStart
            \resumeItem{Custom Autograd Engine}
                {Designed and implemented a custom autograd engine in Scala, enabling automatic computation graph construction and backpropagation.}
            \resumeItem{Gradient Caching for Efficiency}
                {Integrated gradient caching mechanisms to store intermediate values, reducing redundant recomputation and improving training efficiency.}
            \resumeItem{Dynamic Computation Graph}
                {Implemented a dynamic graph structure that supports complex tensor operations, allowing for efficient forward and backward passes in deep learning models.}
            \resumeItem{Scalability and Modularity}
                {Developed reusable tensor and optimization modules, making the engine extensible for future integration with larger ML frameworks, including potential applications in large-scale model training.}
            \resumeItem{RNN and GRU Forecasting}
                {Manually derived backpropagation through time (BPTT) equations for RNNs and GRUs, analyzed gradient flow behavior, and applied gradient clipping to improve training stability.}
            \resumeItem{Implementation in Scala}
                {Implemented RNNs and GRUs from scratch using the derived equations, achieving a SMAPE reduction from 60\% to 16\% on forecasting tasks.}
            \resumeItem{Recurrent Architecture}
                {Designed modular \texttt{RecurrentBase} and \texttt{RNNCell} components for large-scale forecasting tasks.}
        \resumeItemListEnd

    \resumeSubheading
        {Volunteer Researcher}{University of Georgia}
        {Advised by Profs. Fei Dou and Lakshmish Ramaswamy}{Jan 2025 - May 2025}
        \resumeItemListStart
            \resumeItem{Federated Learning for Carbon Footprint Tracking}
                {Implemented \textbf{Federated Learning} using \textbf{Flower} and \textbf{PyTorch} to track \textbf{carbon emissions of deep learning models (MobileNetV3, ResNet20)} across distributed nodes.}
            \resumeItem{Carbon-Aware Model Evaluation}
                {Analyzed trade-offs between \textbf{model accuracy, energy efficiency, and environmental impact} in federated deep learning settings.}
            \resumeItem{Edge AI Experimentation}
                {Conducted experiments on \textbf{Jetson devices} to evaluate federated learning models (MobileNetV3, ResNet20) under low-power, decentralized settings.}
        \resumeItemListEnd

\resumeSubHeadingListEnd

% Skills
\section{Skills}
\resumeSubHeadingListStart
    \item \textbf{Programming Languages}{: Python, Scala, Java, C, SQL}
    \item \textbf{Machine Learning Libraries}{: PyTorch, TensorFlow, Scikit-Learn, OpenCV}
    \item \textbf{Data Science Tools}{: NumPy, Pandas, Matplotlib, Seaborn, Jupyter, scikit-image}
    \item \textbf{ML Techniques \& Concepts}{: Supervised Learning, Deep Learning, Representation Learning, Video Representation Learning, Latent-Space Prediction, Optimization, Backpropagation, Federated Learning, Differential Privacy, Autograd Systems}
\resumeSubHeadingListEnd

% Projects
\section{Projects}
\resumeSubHeadingListStart

    \resumeSubheading
        {Predictive Video Representation Learning with V-JEPA 2}{Research Project}
        {Video Future Selection and Latent Prediction}{2026}
        \resumeItemListStart
            \resumeItem{V-JEPA Integration and Latent Prediction}
                {Built a JEPA-style video future-selection system using a frozen V-JEPA 2 encoder to extract normalized clip embeddings and a compact PyTorch latent predictor trained with a composite MSE and cosine loss to forecast next-step representations.}
            \resumeItem{End-to-End Data and Training Pipeline}
                {Developed MovingMNIST and real-video ingestion pipelines, standardized clips for the video encoder, generated $K$ candidate futures per example, cached extracted embeddings, and implemented configurable AdamW training with validation-based best-state selection.}
            \resumeItem{Scoring, Baselines, and Evaluation}
                {Implemented frozen V-JEPA, learned latent-predictor, semantic, heuristic, and frozen-feature reranking methods; produced candidate rankings, confidence margins, and reproducible CSV/JSON benchmarks that revealed dataset-dependent trade-offs among scorers.}
            \resumeItem{Grounded Interpretability and Demonstration}
                {Created a deterministic explanation layer that converts structured scoring evidence into human-readable commentary, LLM-ready presenter prompts, and visualization-ready summaries for Jupyter demonstrations and research presentations.}
        \resumeItemListEnd

    \resumeSubheading
        {Representation Learning Under Alternative Learning Rules}{University of Georgia}
        {Independent Course Project in Representation Learning}{Jan 2026 - May 2026}
        \resumeItemListStart
            \resumeItem{Learning Rule Implementations}
                {Independently implemented backpropagation, predictive coding, and Mono-Forward under matched multilayer perceptron architectures, conducting controlled experiments and validation-based model selection on MNIST and CIFAR-10.}
            \resumeItem{Representation Analysis}
                {Analyzed learned representations using layer-wise linear probes, t-SNE, and linear centered kernel alignment (CKA); Mono-Forward achieved 98.38\% test accuracy on MNIST, 54.84\% on CIFAR-10, and the strongest probe performance across both datasets.}
            \resumeItem{Representation Similarity}
                {Quantified learning-rule effects on representation geometry: backpropagation and Mono-Forward achieved a mean same-depth CKA of 0.749 on CIFAR-10, compared with 0.411 and 0.425 for pairings involving predictive coding.}
            \resumeItem{Memory Profiling}
                {Profiled peak CPU memory across matched MLPs with 2 to 10 hidden layers, observing lower measured peak memory for Mono-Forward as network depth increased.}
        \resumeItemListEnd

    \resumeSubheading
        {Differentially Private Optimizer Research Project}{University of Georgia}
        {Course Project in Privacy-Preserving Machine Learning (under Prof. Jaewoo Lee)}{Aug 2024 - Dec 2024}
        \resumeItemListStart
            \resumeItem{DP-Lion: Private Optimizer Design}
                {Developed \textbf{DP-Lion}, a novel differentially private optimizer for deep learning, improving training stability and performance over DP-SGD and DP-Adam.}
            \resumeItem{Benchmarking on CIFAR-10}
                {Achieved \textbf{69\% accuracy} on ResNet20 with privacy guarantees, outperforming most existing DP optimizers in empirical evaluations.}
            \resumeItem{Evaluation Against DeepMind Benchmarks}
                {Integrated techniques from DeepMind's private optimization research and compared DP-Lion against their reported baselines.}
            \resumeItem{Research Focus}
                {Used Opacus to privatize DP-Lion and baseline optimizers (DP-SGD, DP-Adam, DP-RMSprop), conducting empirical comparisons as part of a graduate privacy-preserving ML course project.}
            \resumeItem{Code}
                {\href{https://github.com/CS8960-Privacy-Preserving-Data-Analysis/final-project}{GitHub Repository}}
        \resumeItemListEnd

    \resumeSubheading
        {Spotify Song Genre Classification}{University of Georgia}
        {Course Project in Data Mining (under Prof. Fei Dou)}{Aug 2024 - Dec 2024}
        \resumeItemListStart
            \resumeItem{Data Processing \& Feature Engineering}
                {Processed \textbf{114,000 song samples}, extracting insightful \textbf{audio features} (e.g., \textbf{danceability-valence correlation}).}
            \resumeItem{Model Development}
                {Trained and compared \textbf{Decision Trees, Random Forests, SVM, kNN, and Neural Networks}, achieving \textbf{85.6\% accuracy}.}
            \resumeItem{Experimentation \& Metrics}
                {Designed and evaluated experiments using \textbf{precision-recall, ROC-AUC}, and \textbf{confusion matrices} to optimize classification.}
            \resumeItem{Code}
                {\href{https://github.com/Praveen-Rangavajhula/spotify-genre-prediction}{GitHub Repository}}
        \resumeItemListEnd

    \resumeSubheading
        {Diabetes Progression Prediction}{Kaggle Project}
        {Data-Driven Modeling \& Predictive Analytics}{Sept 2024}
        \resumeItemListStart
            \resumeItem{Exploratory Data Analysis (EDA)}
                {Performed \textbf{data cleaning, feature engineering, and transformation (Box-Cox)} to improve regression model performance.}
            \resumeItem{Predictive Modeling}
                {Built and evaluated a \textbf{Linear Regression model from scratch}, benchmarking against \textbf{Scikit-Learn} implementations.}
            \resumeItem{Model Evaluation \& Diagnostics}
                {Used \textbf{residual analysis, Q-Q plots, VIF checks}, and \textbf{homoscedasticity tests} to diagnose and improve model performance.}
            \resumeItem{Regularization \& Feature Selection}
                {Analyzed the impact of \textbf{Lasso, Ridge, and Elastic Net} regularization techniques to mitigate multicollinearity.}
            \resumeItem{Code}
                {\href{https://github.com/Praveen-Rangavajhula/Machine-Learning-Basics}{GitHub Repository}}
        \resumeItemListEnd

    \resumeSubheading
        {Spatio-Temporal Audio Classification}{University of Georgia}
        {Course Project in Machine Learning in IoT (under Prof. Fei Dou)}{Jan 2024 - May 2024}
        \resumeItemListStart
            \resumeItem{Model Adaptation}
                {Adapted a published \textbf{ConvLSTM-based architecture} to classify underwater acoustic signals on the \textbf{VTUAD dataset}.}
            \resumeItem{Learning Outcome}
                {Gained initial hands-on exposure to deep learning pipelines, model reproduction, and working with spatio-temporal data in ML.}
        \resumeItemListEnd

\resumeSubHeadingListEnd

% End
\end{document}